%% =====================================================================
%%  Capítulo 3 -- Modelado II: linealización, segundo orden, bloques y Mason
%%  Migrado al formato nuevo (noviembre de 2026) desde Handout2_CtrlParte2.tex
%%  y las presentaciones Modelado3, Modelado4 y Modelado5.
%%  Cambios de fondo respecto a la versión anterior (ver MIGRACION.md):
%%    el ejemplo numérico de los dos tanques (tau1=1, tau2=4, C2=2) no era
%%    físico (C2 = -C1 siempre, y la planta es SIEMPRE sobreamortiguada): se
%%    usan datos físicos y el subamortiguado se ilustra con masa-resorte;
%%    beta = acos(zeta) (decía cos zeta); MP -> %OS; "determinante" ->
%%    discriminante; "asentamiento" -> establecimiento; los números del
%%    ejemplo con cero y polo adicionales se recalcularon (54.5 % / 38.0 %).
%% =====================================================================

\capitulo[cap:so]{Modelado II: linealización, segundo orden y diagramas de bloques}
  {Mi misión es matar el tiempo, la suya, matarme a mí. Se está perfectamente a gusto entre asesinos.}
  {Emil Michel Cioran}

\begin{objetivos}[partes/03-modelado-segundo-orden.html]
  \item Linealizar un modelo no lineal alrededor de un punto de operación y obtener su función de transferencia en variables desviación.
  \item Reconocer un sistema de segundo orden en forma estándar, clasificarlo por su $\zeta$ y ubicar sus polos.
  \item Calcular $\tr$, $\tp$, $\OS$ y $\ts$ de un subamortiguado, y traducir especificaciones de tiempo en una región del plano-$s$.
  \item Reducir diagramas de bloques con álgebra de bloques y con la ley de Mason, y escribir la ecuación característica.
  \item Decidir cuándo un sistema de orden superior se puede aproximar por sus polos dominantes y simplificarlo, verificando en \matlab/\simulink.
\end{objetivos}

%% ---------------------------------------------------------------------
\section{¿Y si el sistema no es lineal?}
\label{sec:so-nolineal}
%% ---------------------------------------------------------------------

En el \cref{cap:po} todo era lineal: el tanque agitado, con sus suposiciones, daba una ODE lineal de primer orden y la transformada de Laplace hacía el resto. Pero la mayoría de los procesos reales no son lineales: el caudal por una válvula depende de la \emph{raíz cuadrada} de la diferencia de presiones, la velocidad de una reacción del \emph{cuadrado} de una concentración, la fuerza de un resorte endurecido del \emph{cubo} de la elongación. Con un sistema no lineal:
\begin{itemize}
  \item No se puede usar la transformada de Laplace (no es lineal), así que no hay función de transferencia.
  \item El análisis cambia: por lo general se busca entender cualitativamente el comportamiento más que calcularlo exactamente.
\end{itemize}
¿Podemos rescatar algo de la teoría que ya tenemos? Sí: si el proceso opera \emph{cerca} de un punto de operación, se comporta casi como un sistema lineal. Esa aproximación, la \clave{linealización}, es lo primero que haremos. Con ella llegaremos a una función de transferencia de \emph{segundo orden}, el otro gran protagonista de los modelos orientados al control, y a la necesidad de describir sistemas con varios bloques interconectados.

\begin{figure}[ht]
\centering
\includegraphics[width=0.85\linewidth]{cap03/dos_tanques}
\caption{Dos tanques que interactúan. El caudal $f_1$ depende de \emph{las dos} alturas: lo que pasa en un tanque afecta al otro. Adaptado de~\citep{SmCo06}.}
\label{fig:so-tanques}
\end{figure}

\textbf{El problema guía.} En el sistema de la \cref{fig:so-tanques}~\citep{SmCo06}: (1) determinar cómo afecta el caudal de entrada $f_i(t)$ al nivel del segundo tanque $h_2(t)$; (2) escribir el modelo; (3) obtener una función de transferencia entre esa entrada y esa salida; y (4) representar el sistema con un diagrama de bloques. Lo resolvemos en la \cref{sec:so-tanques}, después de aprender a linealizar.

%% ---------------------------------------------------------------------
\section{Linealización}
\label{sec:so-linealizacion}
%% ---------------------------------------------------------------------

La idea es vieja y sencilla: \emph{de cerca, toda curva suave parece una recta}. Si el proceso se mueve poco alrededor de un punto, reemplazamos cada función no lineal por su recta tangente (su serie de Taylor truncada en el primer orden). La suposición de fondo es que la respuesta de la aproximación lineal representa la del proceso en la región cercana al punto de operación~\citep{SmCo06,DoBi08}.

\paragraph{Variables desviación} Como en el \cref{cap:po}, medimos cada variable desde el punto de operación:
\begin{equation}
  \dev{x}(t)=x(t)-\eq{x}.
  \label{eq:so-desviacion}
\end{equation}
Es mover el cero del eje hasta $\eq{x}$; las derivadas no cambian: $\dif^n\dev{x}/\dif t^n=\dif^n x/\dif t^n$. En casi todos los casos el punto de operación será un \clave{punto de equilibrio} (donde las derivadas se anulan); cuando no lo sea, lo diremos.

\paragraph{Una variable} Sea $\dot x=f(x)$ con $f$ suave.\nota{Continuamente diferenciable; para la existencia y unicidad de la solución basta, además, que sea Lipschitz.} Su serie de Taylor alrededor de $x_0$ es
\[
  f(x)=\underbrace{f(x_0)+\left.\deriv{f}{x}\right|_{x_0}(x-x_0)}_{\text{lo que se conserva}}
  +\underbrace{\frac{1}{2!}\left.\frac{\dif^2f}{\dif x^2}\right|_{x_0}(x-x_0)^2+\cdots}_{\text{lo que se desprecia}}
\]
y la aproximación lineal es
\begin{equation}
  f(x)\approx f(x_0)+\left.\deriv{f}{x}\right|_{x_0}\,\dev{x}.
  \label{eq:so-taylor1}
\end{equation}
Si $x_0=\eq{x}$ es un equilibrio, $f(\eq{x})=0$ y queda $\dot{\dev{x}}=a\,\dev{x}$ con $a=f'(\eq{x})$: un modelo lineal en variables desviación.

\begin{ejemplo}[label=ej:so-lin1]{Linealizar alrededor de un punto que no es de equilibrio}
Linealice $\dot x=x\sqrt{x^2+5x}$ alrededor de $x_0=5$.
\solucion La derivada es
\[
  \deriv{f}{x}=\sqrt{x^2+5x}+\frac{x(2x+5)}{2\sqrt{x^2+5x}},
\]
y en $x_0=5$: $f(5)=5\sqrt{50}=35.36$ y $f'(5)=\sqrt{50}+\dfrac{75}{2\sqrt{50}}=12.37$. Entonces
\[
  \dot{\dev{x}}\approx35.36+12.37\,\dev{x}.
\]
Como $x_0=5$ \emph{no} es un equilibrio (los equilibrios son $x=0$ y $x=-5$), queda un término constante: la aproximación es \emph{afín}. La \cref{fig:so-lin1} muestra que solo es buena cerca de $x_0$: el error relativo es menor que el $5\%$ entre $x\approx3.9$ y $x\approx6.7$.
\end{ejemplo}

\begin{figure}[ht]
\centering
\includegraphics[width=0.85\linewidth]{cap03/f_linealizacion}
\caption{Ejemplo~\ref{ej:so-lin1}: la función y su recta tangente en $x_0=5$. En la franja verde el error relativo es menor que el $5\%$.}
\label{fig:so-lin1}
\end{figure}

\codigorepo{cap03_modelado_segundo_orden/C3_Ej1_Linealizacion.m}{simbolica}{ejemplo \ref{ej:so-lin1}: derivada simbólica}

\paragraph{Varias variables} Si $f$ depende de $x_1,\dots,x_n$ (o de estados y entradas), se suma una derivada parcial por variable:
\begin{equation}
  f(x_1,\dots,x_n)\approx f(\eq{x}_1,\dots,\eq{x}_n)+\sum_{i=1}^n\left.\frac{\partial f}{\partial x_i}\right|_{\ast}\dev{x}_i,
  \label{eq:so-taylorn}
\end{equation}
donde ``$\ast$'' indica que las derivadas se evalúan en el punto de operación.

\begin{ejemplo}[label=ej:so-reactor]{Un reactor}
La concentración $x$ de un producto evoluciona según $\dot x=f(x,u)=-kx^2+\tfrac1V\,u\,(c-x)$, donde $u$ es el caudal de alimentación, $k=1$\,L/(mol\,min), $V=1$\,L y $c=2$\,mol/L~\citep{Mar00}. Linealice alrededor de $\eq{x}=1$\,mol/L, $\eq{u}=1$\,L/min.
\solucion Primero, es un equilibrio: $f(1,1)=-1+1\cdot(2-1)=0$. Las derivadas parciales son
\[
  \left.\frac{\partial f}{\partial x}\right|_\ast=-2k\eq{x}-\frac{\eq{u}}{V}=-3,\qquad
  \left.\frac{\partial f}{\partial u}\right|_\ast=\frac{c-\eq{x}}{V}=1,
\]
de modo que $\dot{\dev{x}}=-3\,\dev{x}+\dev{u}$ y, con Laplace,
\[
  \frac{\Dev{X}(s)}{\Dev{U}(s)}=\frac{1}{s+3}=\frac{1/3}{\tfrac13 s+1}:
\]
un primer orden con $K_1=1/3$ y $\tau=1/3$\,min.
\end{ejemplo}

\codigorepo{cap03_modelado_segundo_orden/C3_Ej1_Linealizacion.m}{reactor}{ejemplo \ref{ej:so-reactor}: \cmd{jacobian}}

\begin{receta}[Linealizar un modelo]
\begin{enumerate}
  \item Escribir el modelo como $\dot x_i=f_i(x,u)$ y hallar el punto de operación: resolver $f_i(\eq{x},\eq{u})=0$.
  \item Calcular las derivadas parciales de cada $f_i$ respecto a cada estado y cada entrada, y evaluarlas en el punto.
  \item Escribir el modelo en variables desviación: $\dot{\dev{x}}_i=\sum_j\frac{\partial f_i}{\partial x_j}\big|_\ast\dev{x}_j+\sum_k\frac{\partial f_i}{\partial u_k}\big|_\ast\dev{u}_k$.
  \item Transformar con condiciones iniciales nulas y despejar la función de transferencia.
  \item Verificar contra el modelo no lineal con escalones pequeños \emph{y} grandes: el lineal solo vale cerca del punto.
\end{enumerate}
\end{receta}

\begin{pausa}
Linealice el tanque con descarga libre $A\dot h=q-c\sqrt h$ ($A=2$\,m$^2$, $c=0.5$\,m$^{5/2}$/s) alrededor de $\eq{q}=1$\,m$^3$/s. ¿Cuánto valen $K_1$ y $\tau$?
\respuesta{Equilibrio: $c\sqrt{\eq{h}}=1\Rightarrow\eq{h}=4$\,m. Derivada: $\partial(c\sqrt h)/\partial h=c/(2\sqrt{\eq{h}})=0.125$. Entonces $2\,\dot{\dev{h}}=\dev{q}-0.125\,\dev{h}$ y $\dfrac{\Dev{H}}{\Dev{Q}}=\dfrac{8}{16s+1}$: $K_1=8$\,s/m$^2$ y $\tau=16$\,s.}
\end{pausa}

%% ---------------------------------------------------------------------
\section{Dos tanques que interactúan}
\label{sec:so-tanques}
%% ---------------------------------------------------------------------

\paragraph{Modelo} Los dos tanques de la \cref{fig:so-tanques} están abiertos a la atmósfera y las válvulas descargan a presión atmosférica, así que la diferencia de presión en cada válvula es
\[
  \Delta P_1=(P_A+\rho g h_1)-(P_A+\rho g h_2)=\rho g\,(h_1-h_2),\qquad
  \Delta P_2=(P_A+\rho g h_2)-P_A=\rho g\,h_2,
\]
con $\rho$ la densidad (kg/m$^3$), $g=9.8$\,m/s$^2$ y $h_i$ los niveles (m). Con las aperturas fijas, el caudal por una válvula es $f=C_v\sqrt{\Delta P/G_f}$ ($C_v$: coeficiente de la válvula; $G_f$: gravedad específica), así que
\begin{equation}
  f_1=C'_{v1}\sqrt{h_1-h_2},\qquad f_2=C'_{v2}\sqrt{h_2},\qquad C'_{vi}=C_{vi}\sqrt{\rho g/G_f}.
  \label{eq:so-valvulas}
\end{equation}
El caudal $f_1$ depende de los dos niveles: lo que pasa en un tanque se refleja en el otro y viceversa. El sistema está \emph{acoplado}. Los balances de masa (con densidad constante y áreas $A_1$, $A_2$) son
\begin{equation}
  A_1\deriv{h_1}{t}=f_i-C'_{v1}\sqrt{h_1-h_2},\qquad
  A_2\deriv{h_2}{t}=C'_{v1}\sqrt{h_1-h_2}-C'_{v2}\sqrt{h_2}.
  \label{eq:so-tanques-nl}
\end{equation}
Son no lineales: hay que linealizar.

\paragraph{Equilibrio} Operamos con $\eq{f_i}=\eq{f_1}=\eq{f_2}$, que fija los niveles $\eq{h_1}$ y $\eq{h_2}$ (las derivadas se anulan).

\paragraph{Linealización} Aplicando \eqref{eq:so-taylorn} a \eqref{eq:so-valvulas},
\[
  f_1\approx\eq{f_1}+C_1\,\dev{h}_1+C_2\,\dev{h}_2,\qquad f_2\approx\eq{f_2}+C_3\,\dev{h}_2,
\]
\begin{equation}
  C_1=\left.\frac{\partial f_1}{\partial h_1}\right|_\ast=\frac{C'_{v1}}{2\sqrt{\eq{h_1}-\eq{h_2}}},\qquad
  C_2=\left.\frac{\partial f_1}{\partial h_2}\right|_\ast=-C_1,\qquad
  C_3=\left.\frac{\partial f_2}{\partial h_2}\right|_\ast=\frac{C'_{v2}}{2\sqrt{\eq{h_2}}}.
  \label{eq:so-coeficientes}
\end{equation}
Reemplazando en los balances y restando el equilibrio,
\[
  A_1\,\dot{\dev{h}}_1=\dev{f}_i-C_1\dev{h}_1-C_2\dev{h}_2,\qquad
  A_2\,\dot{\dev{h}}_2=C_1\dev{h}_1+(C_2-C_3)\,\dev{h}_2.
\]

\paragraph{Función de transferencia} Transformando y despejando cada nivel,
\[
  \Dev{H}_1(s)=\frac{K_1}{\tau_1s+1}\bigl(\Dev{F}_i(s)-C_2\Dev{H}_2(s)\bigr),\qquad
  \Dev{H}_2(s)=\frac{K_2}{\tau_2s+1}\,\Dev{H}_1(s),
\]
con $K_1=1/C_1$, $\tau_1=A_1/C_1$, $K_2=C_1/(C_3-C_2)$ y $\tau_2=A_2/(C_3-C_2)$. Eliminando $\Dev{H}_1$:
\begin{equation}
  \frac{\Dev{H}_2(s)}{\Dev{F}_i(s)}=\frac{K_1K_2}{\tau_1\tau_2s^2+(\tau_1+\tau_2)s+1+K_1K_2C_2}.
  \label{eq:so-tanques-ft}
\end{equation}
Las dos ecuaciones de primer orden, acopladas, producen un sistema de \clave{segundo orden}. (La estructura de bloques de esta deducción está en el problema~\ref{prob:so-bloques-tanques}.)

\begin{ejemplo}[label=ej:so-tanques]{Los tanques con números}
Con $A_1=2$\,m$^2$, $A_2=3$\,m$^2$, $C'_{v1}=0.1$ y $C'_{v2}=0.05$\,m$^{5/2}$/s, y el tanque 2 operando en $\eq{h_2}=4$\,m: halle el equilibrio, la función de transferencia y sus polos.
\solucion Equilibrio: $\eq{f}=C'_{v2}\sqrt{4}=0.1$\,m$^3$/s y $\eq{h_1}-\eq{h_2}=(0.1/0.1)^2=1$, luego $\eq{h_1}=5$\,m. Coeficientes: $C_1=0.1/(2\cdot1)=0.05$, $C_2=-0.05$, $C_3=0.05/(2\cdot2)=0.0125$\,m$^2$/s. Entonces $K_1=20$, $\tau_1=40$\,s, $K_2=0.8$, $\tau_2=48$\,s y
\[
  \frac{\Dev{H}_2(s)}{\Dev{F}_i(s)}=\frac{16}{1920s^2+88s+0.2}=\frac{80}{9600s^2+440s+1}.
\]
La ganancia es $80$\,s/m$^2$ ($=1/C_3$: en estado estacionario todo lo que entra sale por la segunda válvula). Los polos son $-0.0024$ y $-0.0434$\,s$^{-1}$: \emph{reales}. La \cref{fig:so-tanques-nl} compara el modelo lineal con el no lineal: con un escalón de $0.01$\,m$^3$/s ($10\%$ del caudal) el error final es de $4$\,cm; con $0.05$\,m$^3$/s, de casi un metro.
\end{ejemplo}

\begin{figure}[ht]
\centering
\includegraphics[width=\linewidth]{cap03/f_tanques}
\caption{Ejemplo~\ref{ej:so-tanques}: nivel del tanque 2 con el modelo no lineal \eqref{eq:so-tanques-nl} y con el linealizado. Cuanto más grande el escalón, más nos alejamos del punto de operación y peor la aproximación.}
\label{fig:so-tanques-nl}
\end{figure}

\codigorepo{cap03_modelado_segundo_orden/C3_Ej2_Tanques.m}{tanques}{ejemplo \ref{ej:so-tanques}}

\begin{cuidado}[Esta planta nunca oscila]
Como $C_2=-C_1$ \emph{siempre} (lo que sube $f_1$ al subir $h_1$ es lo que baja al subir $h_2$), la ecuación característica se puede escribir $s^2+\bigl(\tfrac{C_1+C_3}{A_2}+\tfrac{C_1}{A_1}\bigr)s+\tfrac{C_1C_3}{A_1A_2}$, y la desigualdad entre las medias aritmética y geométrica muestra que $\zeta\geq1$ para cualesquiera parámetros físicos (problema~\ref{prob:so-zeta-tanques}): \alerta{dos tanques que interactúan por gravedad nunca dan una respuesta subamortiguada}. La versión anterior de estas notas usaba $\tau_1=1$, $\tau_2=4$, $C_2=2$, $K_1=K_2=1$ para mostrar un sobrepico de $3.78\%$; esos valores no corresponden a ningún par de tanques reales ($C_2>0$ exigiría que el caudal entre tanques \emph{aumentara} al llenarse el segundo). Para ver oscilaciones usaremos un sistema mecánico (\cref{sec:so-segundoorden}).
\end{cuidado}

\begin{modelo}[Tanques.slx]
El script \cmd{crearModelo\_Tanques.m} calcula el equilibrio y la función de transferencia, y construye por código el modelo no lineal (integradores y bloques \emph{Fcn} con las raíces cuadradas) al lado del linealizado, con el mismo escalón en $f_i$. Pruebe \cmd{df = 0.01} y \cmd{df = 0.05} y cuantifique el error final.
\end{modelo}

%% ---------------------------------------------------------------------
\section{El sistema de segundo orden}
\label{sec:so-segundoorden}
%% ---------------------------------------------------------------------

La función \eqref{eq:so-tanques-ft} y muchas otras se pueden escribir en la \clave{forma estándar del segundo orden}:
\begin{equation}
  T(s)=\frac{Y(s)}{U(s)}=K\,\frac{\wn^2}{s^2+2\zeta\wn s+\wn^2},
  \label{eq:so-estandar}
\end{equation}
con tres parámetros: la \clave{ganancia} $K$ (el valor final ante un escalón unitario), la \clave{frecuencia natural} $\wn$ (rad/s) y el \clave{factor de amortiguamiento} $\zeta$ (adimensional). Para los tanques, $\wn^2=\frac{1+K_1K_2C_2}{\tau_1\tau_2}$ y $2\zeta\wn=\frac{\tau_1+\tau_2}{\tau_1\tau_2}$. Los polos son
\begin{equation}
  p_{1,2}=-\zeta\wn\pm\wn\sqrt{\zeta^2-1},
  \label{eq:so-polos}
\end{equation}
y su naturaleza depende solo de $\zeta$:
\begin{center}
\begin{tabular}{@{}lll@{}}
\toprule
$\zeta$ & polos & respuesta al escalón\\
\midrule
$\zeta=0$ & imaginarios puros $\pm j\wn$ & oscilación sostenida (\emph{no amortiguado})\\
$0<\zeta<1$ & complejos: $-\zeta\wn\pm j\wdamp$ & oscilación que se extingue (\clave{subamortiguado})\\
$\zeta=1$ & reales e iguales: $-\wn$ (doble) & la más rápida sin sobrepico (\clave{críticamente amortiguado})\\
$\zeta>1$ & reales y distintos & lenta, sin sobrepico (\clave{sobreamortiguado})\\
\bottomrule
\end{tabular}
\end{center}
donde $\wdamp=\wn\sqrt{1-\zeta^2}$ es la \clave{frecuencia natural amortiguada}. La \cref{fig:so-zetas} lo resume.

\begin{figure}[ht]
\centering
\includegraphics[width=\linewidth]{cap03/f_zetas}
\caption{Segundo orden con $\wn=1$. Izquierda: al aumentar $\zeta$ los polos recorren la circunferencia de radio $\wn$ hasta juntarse en $-1$ ($\zeta=1$) y luego se separan sobre el eje real. Derecha: las respuestas al escalón correspondientes.}
\label{fig:so-zetas}
\end{figure}

\codigorepo{cap03_modelado_segundo_orden/C3_Ej3_SegundoOrden.m}{zetas}{\cref{fig:so-zetas}}

\subsection{Un segundo orden que sí oscila: masa, resorte y amortiguador}

La \cref{fig:so-masa} es el prototipo mecánico del segundo orden (una suspensión, un edificio en un sismo, el brazo de un disco duro). La segunda ley de Newton da
\begin{equation}
  m\ddot x+b\dot x+kx=F
  \quad\Longrightarrow\quad
  \frac{X(s)}{F(s)}=\frac{1}{ms^2+bs+k}=\frac1k\,\frac{k/m}{s^2+\frac bm s+\frac km},
  \label{eq:so-masa}
\end{equation}
así que $K=1/k$, $\wn=\sqrt{k/m}$ y $\zeta=\dfrac{b}{2\sqrt{km}}$. Con poco amortiguamiento ($b$ pequeño) la masa oscila: la energía va y viene entre el resorte y la masa, y el amortiguador la va disipando. Con $m=1$\,kg, $b=1.2$\,N\,s/m y $k=4$\,N/m resultan $\wn=2$\,rad/s y $\zeta=0.3$.

\begin{figure}[ht]
\centering
\begin{minipage}{0.4\linewidth}\centering
\includegraphics[width=\linewidth]{cap03/masa_resorte}
\end{minipage}\hfill
\begin{minipage}{0.57\linewidth}\centering
\includegraphics[width=\linewidth]{cap03/f_masa_resorte}
\end{minipage}
\caption{Masa-resorte-amortiguador y su respuesta a un escalón de $1$\,N. El valor final es $1/k=0.25$\,m y el sobrepico, $37\%$.}
\label{fig:so-masa}
\end{figure}

\begin{pausa}
En la masa-resorte de la \cref{fig:so-masa}, ¿qué valor de $b$ la vuelve críticamente amortiguada?
\respuesta{$\zeta=1\Rightarrow b=2\sqrt{km}=2\sqrt{4}=4$\,N\,s/m.}
\end{pausa}

%% ---------------------------------------------------------------------
\section{La respuesta del subamortiguado, en números}
\label{sec:so-caracteristicas}
%% ---------------------------------------------------------------------

El escalón es la entrada de prueba clásica: es un cambio brusco, fácil de generar y suficientemente drástico para ver la dinámica (y, conocida su respuesta, se obtienen las demás por linealidad). Tomemos $K=1$, condiciones iniciales nulas y $0<\zeta<1$. Con $U(s)=1/s$,
\[
  Y(s)=\frac{\wn^2}{s(s^2+2\zeta\wn s+\wn^2)}
  =\frac1s-\frac{s+\zeta\wn}{(s+\zeta\wn)^2+\wdamp^2}-\frac{\zeta\wn}{(s+\zeta\wn)^2+\wdamp^2},
\]
y con las dos últimas filas de la tabla de transformadas,
\[
  y(t)=1-e^{-\zeta\wn t}\left(\cos\wdamp t+\frac{\zeta}{\sqrt{1-\zeta^2}}\sin\wdamp t\right).
\]
Con la identidad $C\cos\theta+D\sin\theta=\sqrt{C^2+D^2}\,\sin\bigl(\theta+\tan^{-1}(C/D)\bigr)$ queda
\begin{equation}
  y(t)=1-\frac{e^{-\zeta\wn t}}{\sqrt{1-\zeta^2}}\,\sin\left(\wdamp t+\beta\right),\qquad \beta=\cos^{-1}\zeta.
  \label{eq:so-subamortiguado}
\end{equation}
La salida oscila a la frecuencia $\wdamp$ dentro de las \clave{envolventes} $1\pm e^{-\zeta\wn t}/\sqrt{1-\zeta^2}$, y el error $1-y(t)$ tiende a cero. (Con $\zeta=0$ no: oscila para siempre a $\wn$.)

El ángulo $\beta$ tiene una lectura geométrica muy útil (\cref{fig:so-beta}): es el ángulo entre el semieje real negativo y la recta que une el origen con el polo. Como la distancia del polo al origen es $\wn$ y su parte real es $-\zeta\wn$, $\cos\beta=\zeta$.

\begin{figure}[ht]
\centering
\includegraphics[width=0.55\linewidth]{cap03/plano_s_beta}
\caption{Geometría de los polos subamortiguados: distancia al origen $\wn$, parte real $-\zeta\wn$, parte imaginaria $\wdamp$ y $\cos\beta=\zeta$.}
\label{fig:so-beta}
\end{figure}

\subsection{Las cuatro características}

\begin{figure}[ht]
\centering
\includegraphics[width=0.95\linewidth]{cap03/f_caracteristicas}
\caption{Características de la respuesta al escalón con $\zeta=0.4$ y $\wn=1$. La franja verde es la banda del $\pm2\%$; las líneas punteadas, las envolventes.}
\label{fig:so-caracteristicas}
\end{figure}

\paragraph{Tiempo de subida $\tr$} Hay varias definiciones: del $10$ al $90\%$ del valor final (la de \cmd{stepinfo} y la usual para sobreamortiguados), del $5$ al $95\%$, o del $0$ al $100\%$ (común para subamortiguados). Para la de $0$--$100\%$, $y(\tr)=1$ en \eqref{eq:so-subamortiguado} exige $\sin(\wdamp\tr+\beta)=0$, es decir $\wdamp\tr+\beta=\pi$:
\begin{equation}
  \tr=\frac{\pi-\beta}{\wdamp}\quad(0\text{--}100\%).
  \label{eq:so-tr}
\end{equation}
Para la de $10$--$90\%$ no hay fórmula cerrada; un ajuste a la solución exacta es
\begin{equation}
  \tr\approx\frac{2.3\zeta^2-0.078\zeta+1.12}{\wn}\quad(10\text{--}90\%),
  \label{eq:so-tr1090}
\end{equation}
con error menor que el $3.1\%$ para $0.1\leq\zeta\leq0.9$ (lo comprobamos numéricamente; otra aproximación muy citada, $(1-0.4167\zeta+2.917\zeta^2)/\wn$, se equivoca hasta un $14\%$ con $\zeta$ pequeño).

\paragraph{Tiempo de pico $\tp$} Es el instante del primer máximo. Derivando \eqref{eq:so-subamortiguado} e igualando a cero se llega a $\frac{\wn}{\sqrt{1-\zeta^2}}e^{-\zeta\wn t}\sin\wdamp t=0$; el primer cero positivo es $\wdamp\tp=\pi$:
\begin{equation}
  \tp=\frac{\pi}{\wdamp}.
  \label{eq:so-tp}
\end{equation}

\paragraph{Sobrepico porcentual $\OS$} Es cuánto se pasa el máximo del valor final, en porcentaje: $\OS=\frac{y(\tp)-y(\infty)}{y(\infty)}\times100\%$. Evaluando en $\tp$:
\begin{equation}
  \OS=100\,e^{-\zeta\pi/\sqrt{1-\zeta^2}}\ \%,
  \qquad\text{y al revés,}\qquad
  \zeta=\frac{-\ln(\OS/100)}{\sqrt{\pi^2+\ln^2(\OS/100)}}.
  \label{eq:so-os}
\end{equation}
\emph{Solo depende de $\zeta$} (\cref{fig:so-os}). Más $\zeta$, menos sobrepico, pero respuesta más lenta.

\begin{figure}[ht]
\centering
\includegraphics[width=0.8\linewidth]{cap03/f_os_zeta}
\caption{Sobrepico en función de $\zeta$, con algunos valores que aparecen en el curso ($\zeta=0.4559$ da $20\%$; $\zeta=0.5912$, $10\%$; $\zeta=0.707$, $4.3\%$).}
\label{fig:so-os}
\end{figure}

\paragraph{Tiempo de establecimiento $\ts$} Es el tiempo a partir del cual la respuesta se queda dentro de una banda alrededor del valor final ($\pm2\%$ en el curso; a veces $\pm5\%$). Usando las envolventes, que decaen con constante de tiempo $1/(\zeta\wn)$, y el mismo criterio del primer orden ($4\tau$ para el $2\%$, $3\tau$ para el $5\%$):
\begin{equation}
  \ts\approx\frac{4}{\zeta\wn}\ (2\%),\qquad \ts\approx\frac{3}{\zeta\wn}\ (5\%).
  \label{eq:so-ts}
\end{equation}
Es una estimación por la envolvente: el valor exacto puede ser menor (en la \cref{fig:so-caracteristicas}, $8.41$ en lugar de $10$), porque la oscilación entra a la banda antes que la envolvente.

\begin{table}[ht]
\centering
\caption{Resumen del segundo orden subamortiguado $\wn^2/(s^2+2\zeta\wn s+\wn^2)$.}
\label{tab:so-resumen}
\renewcommand{\arraystretch}{1.9}
\begin{tabular}{@{}lll@{}}
\toprule
característica & fórmula & depende de\\
\midrule
tiempo de subida (0--100\%) & $\tr=\dfrac{\pi-\cos^{-1}\zeta}{\wn\sqrt{1-\zeta^2}}$ & $\zeta$ y $\wn$\\
tiempo de pico & $\tp=\dfrac{\pi}{\wn\sqrt{1-\zeta^2}}$ & parte imaginaria del polo\\
sobrepico & $\OS=100\,e^{-\zeta\pi/\sqrt{1-\zeta^2}}$ & solo $\zeta$ (ángulo del polo)\\
establecimiento (2\%) & $\ts\approx\dfrac{4}{\zeta\wn}$ & parte real del polo\\
\bottomrule
\end{tabular}
\end{table}

\codigorepo{cap03_modelado_segundo_orden/C3_Ej3_SegundoOrden.m}{formulas}{fórmulas contra \cmd{stepinfo}}

\begin{cuidado}[¿Qué $\tr$ reporta \matlab?]
\cmd{stepinfo} usa $10$--$90\%$. Si compara con \eqref{eq:so-tr} (de $0$--$100\%$) le van a dar números distintos y ninguno está mal. Diga siempre qué definición usa.
\end{cuidado}

\subsection{Moverse en el plano-$s$}

La \cref{tab:so-resumen} dice algo que conviene ver geométricamente: con los polos en $-\zeta\wn\pm j\wdamp$,
\begin{itemize}
  \item sobre una \emph{recta vertical} (parte real fija) el $\ts$ es constante;
  \item sobre una \emph{recta horizontal} (parte imaginaria fija) el $\tp$ es constante;
  \item sobre una \emph{recta radial} (ángulo $\beta$ fijo) el $\zeta$, y con él el $\OS$, son constantes.
\end{itemize}
La \cref{fig:so-lineas} muestra qué pasa al moverse por cada una. Por eso las especificaciones ``$\OS$ máximo y $\ts$ máximo'' se convierten en una \emph{región} del plano-$s$ donde deben quedar los polos: a la izquierda de una vertical y dentro de un cono alrededor del semieje real negativo. Esa idea es la base del lugar de las raíces.

\begin{figure}[ht]
\centering
\includegraphics[width=\linewidth]{cap03/f_lineas_plano}
\caption{Moviendo los polos por rectas del plano-$s$. Vertical: mismo $\ts$, más oscilación hacia arriba. Horizontal: mismo $\tp$, menos sobrepico hacia la izquierda. Radial: mismo $\OS$, más rápido lejos del origen.}
\label{fig:so-lineas}
\end{figure}

\subsection{Críticamente amortiguado y sobreamortiguado}

\paragraph{$\zeta=1$} Hay un polo doble en $-\wn$ y
\begin{equation}
  y(t)=1-e^{-\wn t}\,(1+\wn t),
  \label{eq:so-critico}
\end{equation}
que también se obtiene de \eqref{eq:so-subamortiguado} con $\lim_{\zeta\to1}\sin(\wdamp t)/\sqrt{1-\zeta^2}=\wn t$. Es la respuesta más rápida \emph{sin} sobrepico.

\paragraph{$\zeta>1$} Hay dos polos reales en $-s_1$ y $-s_2$, con $s_{1,2}=\zeta\wn\pm\wn\sqrt{\zeta^2-1}$, y
\begin{equation}
  y(t)=1+\frac{\wn}{2\sqrt{\zeta^2-1}}\left(\frac{e^{-s_1t}}{s_1}-\frac{e^{-s_2t}}{s_2}\right).
  \label{eq:so-sobre}
\end{equation}
Si $s_1\gg s_2$, el término en $e^{-s_1t}$ muere muy rápido y la respuesta se parece a la de un primer orden con el polo lento: $T(s)\approx\dfrac{s_2}{s+s_2}$. Es la primera aparición de la idea de \emph{polo dominante} (\cref{sec:so-ordensuperior}). Los tanques del ejemplo~\ref{ej:so-tanques} son así: $s_2=0.0024$ y $s_1=0.0434$, y $4/s_2\approx1670$\,s es una buena estimación del $\ts$ exacto ($1655$\,s).

\begin{pausa}
Un sistema tiene polos en $-2\pm j4$. Sin simular: ¿cuánto valen $\wn$, $\zeta$, $\tp$, $\OS$ y $\ts$?
\respuesta{$\wn=\sqrt{4+16}=4.47$\,rad/s, $\zeta=2/4.47=0.447$, $\tp=\pi/4=0.785$\,s, $\OS=100e^{-0.447\pi/0.894}=20.8\%$ y $\ts\approx4/2=2$\,s.}
\end{pausa}

%% ---------------------------------------------------------------------
\section{Diagramas de bloques}
\label{sec:so-bloques}
%% ---------------------------------------------------------------------

Un proceso real tiene válvulas, sensores, transmisores y controladores, cada uno con su propia función de transferencia. Para representar esa cadena de causas y efectos usamos \clave{diagramas de bloques}: elementos \emph{unidireccionales}, cada uno con la función de transferencia entre su entrada y su salida (\cref{fig:so-elementos}).

\begin{figure}[ht]
\centering
\includegraphics[width=\linewidth]{cap03/bloques_elementos}
\caption{Los elementos de un diagrama de bloques: (a) una señal, (b) un comparador, (c) un punto de bifurcación (la señal se copia, no se reparte), (d) un sumador con varios signos y (e) un bloque.}
\label{fig:so-elementos}
\end{figure}

Tres reglas bastan para reducir casi cualquier diagrama (\cref{fig:so-reglas}):
\begin{itemize}
  \item \textbf{Serie:} $Z=G_1R$ y $Y=G_2Z$ dan $Y=G_1G_2\,R$. Las funciones de transferencia en cascada se multiplican.
  \item \textbf{Paralelo:} $Y=G_1R+G_2R=(G_1+G_2)\,R$.
  \item \textbf{Retroalimentación negativa:} con $E=R-\Hs Y$ y $Y=GE$,
  \begin{equation}
    \frac{Y(s)}{R(s)}=\frac{G(s)}{1+G(s)\Hs(s)}.
    \label{eq:so-retro}
  \end{equation}
  Con retroalimentación positiva ($E=R+\Hs Y$) el signo del denominador cambia: $G/(1-G\Hs)$.
\end{itemize}

\begin{figure}[ht]
\centering
\includegraphics[width=0.95\linewidth]{cap03/bloques_reglas}
\caption{Las tres reglas del álgebra de bloques.}
\label{fig:so-reglas}
\end{figure}

\begin{ideaclave}
El denominador de \eqref{eq:so-retro} igualado a cero, $1+G(s)\Hs(s)=0$, es la \clave{ecuación característica} del lazo cerrado. El producto $G\Hs$ es la \clave{malla abierta} $\Gma(s)$ (\cref{def:rf-malla}): lo que recorre una señal al dar una vuelta completa al lazo. Casi todo el análisis del curso (estabilidad, lugar de las raíces, Bode) gira alrededor de esa función.
\end{ideaclave}

\begin{ejemplo}[label=ej:so-dosentradas]{Dos entradas, dos salidas}
Para el diagrama de la \cref{fig:so-dosentradas}, demuestre que $Y_2(s)=G_3(s)\bigl[G_1(s)R_1(s)+G_2(s)R_2(s)\bigr]$.
\solucion La salida del sumador es $G_1R_1+G_2R_2$; el punto de bifurcación la copia en $Y_1$ y en la entrada de $G_3$. Por la regla de la serie, $Y_2=G_3Y_1=G_3(G_1R_1+G_2R_2)$. Por linealidad, cada entrada contribuye por su propio camino (superposición).
\end{ejemplo}

\begin{figure}[ht]
\centering
\includegraphics[width=0.75\linewidth]{cap03/bloques_dos_entradas}
\caption{Ejemplo~\ref{ej:so-dosentradas}.}
\label{fig:so-dosentradas}
\end{figure}

No siempre se quiere \emph{simplificar}: a veces se busca que el diagrama tenga cierta \emph{estructura}. El caso más importante es reescribir un lazo con retroalimentación $\Hs(s)$ como uno con retroalimentación unitaria, que es la forma en la que se estudia el error en estado estacionario.

\begin{ejemplo}[label=ej:so-equivalente]{Lazo abierto equivalente}
Demuestre que los dos lazos de la \cref{fig:so-equivalente} son equivalentes si
\begin{equation}
  \Gla(s)=\frac{G(s)}{1+G(s)\Hs(s)-G(s)}.
  \label{eq:so-gla}
\end{equation}
\solucion El lazo (b) tiene $T=\Gla/(1+\Gla)$. Con \eqref{eq:so-gla},
\[
  1+\Gla=\frac{1+G\Hs-G+G}{1+G\Hs-G}=\frac{1+G\Hs}{1+G\Hs-G}
  \quad\Longrightarrow\quad
  \frac{\Gla}{1+\Gla}=\frac{G}{1+G\Hs},
\]
que es el lazo cerrado de (a). A $\Gla$ la llamamos \clave{lazo abierto} (\cref{def:rf-malla}); no hay que confundirla con la malla abierta $G\Hs$.
\end{ejemplo}

\begin{figure}[ht]
\centering
\includegraphics[width=0.95\linewidth]{cap03/lazo_equivalente}
\caption{Ejemplo~\ref{ej:so-equivalente}: el mismo lazo cerrado con retroalimentación $\Hs$ y con retroalimentación unitaria.}
\label{fig:so-equivalente}
\end{figure}

En \matlab las tres reglas son \cmd{series}, \cmd{parallel} y \cmd{feedback} (o, simplemente, \cmd{*}, \cmd{+} y \cmd{feedback}); \cmd{minreal} cancela los polos y ceros repetidos que deja el álgebra:
\codigorepo{cap03_modelado_segundo_orden/C3_Ej5_Bloques.m}{algebra}{un lazo interno dentro de uno externo (problema \ref{prob:so-multilazo})}
\codigorepo{cap03_modelado_segundo_orden/C3_Ej5_Bloques.m}{equivalente}{ejemplo \ref{ej:so-equivalente}}

\begin{cuidado}[\texttt{feedback} supone signo negativo]
\cmd{feedback(G,H)} calcula $G/(1+GH)$. Para retroalimentación positiva hay que escribir \cmd{feedback(G,H,+1)}. Y \cmd{feedback(G,1)} es retroalimentación \emph{unitaria}, no ``sin retroalimentación''.
\end{cuidado}

\begin{pausa}
Si en la regla de retroalimentación $G\Hs\gg1$ en las frecuencias de interés, ¿a qué se parece $Y/R$? ¿Qué consecuencia práctica tiene?
\respuesta{$Y/R\approx1/\Hs$: el lazo cerrado depende casi solo del sensor y no de la planta. Por eso un buen sensor es más importante que una planta precisa, y por eso Black pudo construir amplificadores precisos con tubos imprecisos (\cref{cap:rf}).}
\end{pausa}

%% ---------------------------------------------------------------------
\section{Diseñar con especificaciones de tiempo}
\label{sec:so-diseno}
%% ---------------------------------------------------------------------

Las fórmulas de la \cref{sec:so-caracteristicas} sirven al revés: dadas unas especificaciones, encontrar los parámetros que las cumplen. Es nuestro primer diseño.

\begin{ejemplo}[label=ej:so-diseno]{Servo de posición con tacómetro}
En el servo de la \cref{fig:so-kkh}, la posición se mide con un sensor y la velocidad con un tacómetro de ganancia $K_h$. Halle $K$ y $K_h$ para que la respuesta al escalón tenga $\OS=20\%$ y $\tp=1$\,s.
\solucion \textbf{1. Lazo cerrado.} El camino directo es $G=\frac{K}{s}\cdot\frac{1}{s+1}=\frac{K}{s(s+1)}$ y la retroalimentación $\Hs=1+K_hs$. Por \eqref{eq:so-retro},
\[
  \frac{Y(s)}{R(s)}=\frac{K}{s^2+(1+KK_h)\,s+K}.
\]
Es la forma estándar \eqref{eq:so-estandar} con ganancia $1$, $\wn^2=K$ y $2\zeta\wn=1+KK_h$.

\textbf{2. De las especificaciones a $\zeta$ y $\wn$.} Con \eqref{eq:so-os}, $\OS=20\%$ da $\zeta=0.4559$; con \eqref{eq:so-tp}, $\wn=\dfrac{\pi}{\tp\sqrt{1-\zeta^2}}=3.5299$\,rad/s.

\textbf{3. Igualar coeficientes.} $K=\wn^2=12.46$ y $K_h=\dfrac{2\zeta\wn-1}{K}=0.178$\,s. El lazo cerrado queda
\[
  T(s)=\frac{12.46}{s^2+3.219\,s+12.46}.
\]
\textbf{4. Verificar} (\cref{fig:so-diseno}): \cmd{stepinfo} reporta $20.0\%$ y $1.00$\,s.
\end{ejemplo}

\begin{figure}[ht]
\centering
\includegraphics[width=0.85\linewidth]{cap03/lazo_kkh}
\caption{Servo de posición con retroalimentación de posición y de velocidad (tacómetro).}
\label{fig:so-kkh}
\end{figure}

\begin{figure}[ht]
\centering
\includegraphics[width=0.85\linewidth]{cap03/f_diseno_kkh}
\caption{Ejemplo~\ref{ej:so-diseno}. Con $K=12.46$ y sin tacómetro el sobrepico sería del $64\%$: la retroalimentación de velocidad agrega amortiguamiento sin cambiar $\wn$.}
\label{fig:so-diseno}
\end{figure}

\codigorepo{cap03_modelado_segundo_orden/C3_Ej4_DisenoKKh.m}{diseno}{ejemplo \ref{ej:so-diseno}}

\begin{modelo}[Servo.slx]
El script \cmd{crearModelo\_Servo.m} arma el servo como $K\to\frac{1}{s+1}\to\frac1s$: la salida del motor es la velocidad, que el tacómetro realimenta sin necesidad de un bloque derivador. Cambie $K$ y $K_h$ por separado y observe qué hace cada uno con el sobrepico y el tiempo de pico.
\end{modelo}

\begin{pausa}
En el ejemplo~\ref{ej:so-diseno}, ¿qué pasa con el $\OS$ y el $\tp$ si se duplica $K$ dejando $K_h=0.178$?
\respuesta{$\wn=\sqrt{24.92}=4.99$ y $2\zeta\wn=1+24.92\cdot0.178=5.44$, luego $\zeta=0.545$. El sobrepico baja (a $13\%$) y el $\tp$ también ($\pi/(4.99\sqrt{1-0.297})=0.75$\,s): con tacómetro, subir $K$ no necesariamente empeora el sobrepico.}
\end{pausa}

%% ---------------------------------------------------------------------
\section{La ley de Mason}
\label{sec:so-mason}
%% ---------------------------------------------------------------------

Con muchos lazos entrelazados, el álgebra de bloques se vuelve un juego de paciencia. Hay una alternativa que da la respuesta de una vez.

\begin{historia}[Samuel J.\ Mason y los grafos de flujo de señal]
En 1953, Samuel J.\ Mason, del MIT, propuso representar un sistema de ecuaciones lineales como un \emph{grafo de flujo de señal} y mostró cómo leer la retroalimentación en él~\citep{Mas53}; en 1956 publicó la fórmula general de la ganancia entre dos nodos, que tiene en cuenta los lazos que no se tocan~\citep{Mas56}. Su fórmula completó lo que Arnold Tustin había dejado sin resolver en un trabajo previo sobre diagramas de flujo. Hoy se la conoce como \emph{Mason's gain formula} o \clave{ley de Mason}.
\end{historia}

En un grafo de flujo, las señales son \emph{nodos} y cada \emph{rama} dirigida lleva una ganancia (una función de transferencia). Un nodo suma las señales de todas las ramas que le llegan y transmite esa suma por todas las que salen. Vocabulario~\citep{ogata10}:

\begin{definicion}[label=def:so-grafo]{Grafo de flujo de señal}
\begin{itemize}
  \item \textbf{Nodo de entrada} (\emph{source}): solo tiene ramas salientes. \textbf{Nodo de salida} (\emph{sink}): solo entrantes. \textbf{Nodo mixto}: ambas.
  \item \textbf{Trayecto}: un recorrido por ramas en el sentido de las flechas. \textbf{Trayecto directo}: va del nodo de entrada al de salida sin pasar dos veces por el mismo nodo.
  \item \textbf{Lazo}: un trayecto cerrado que no repite nodos. Dos lazos son \textbf{disyuntos} si no tienen ningún nodo en común (no se tocan).
  \item \textbf{Ganancia} de un trayecto o de un lazo: el producto de las ganancias de sus ramas.
\end{itemize}
\end{definicion}

\begin{definicion}[label=def:so-mason]{Ley de Mason}
La función de transferencia entre el nodo de entrada y uno de salida es
\begin{equation}
  T(s)=\frac{1}{\Delta}\sum_{k=1}^{n}P_k\,\Delta_k,
  \label{eq:so-mason}
\end{equation}
donde $n$ es el número de trayectos directos, $P_k$ la ganancia del $k$-ésimo, $\Delta$ el \clave{discriminante} del grafo,
\begin{equation}
  \Delta=1-\sum_a L_a+\sum_{b,c}L_bL_c-\sum_{d,e,f}L_dL_eL_f+\cdots
  \label{eq:so-discriminante}
\end{equation}
(uno, menos la suma de las ganancias de todos los lazos, más la suma de los productos de pares de lazos disyuntos, menos la de ternas de lazos disyuntos, y así sucesivamente), y $\Delta_k$ el discriminante de la parte del grafo que \emph{no toca} el trayecto $P_k$ (se calcula como $\Delta$, eliminando los lazos que tocan $P_k$).
\end{definicion}

\begin{cuidado}[Discriminante, no determinante]
En la versión anterior se llamaba ``determinante'' a $\Delta$. Lo llamamos \clave{discriminante} para no confundirlo con el determinante de una matriz (aunque, en el fondo, es el determinante del sistema de ecuaciones de los nodos, y por eso funciona).
\end{cuidado}

Antes de un caso grande, comprobemos con el lazo de siempre: un trayecto $P_1=G$, un lazo $L_1=-G\Hs$ que toca al trayecto, así que $\Delta=1+G\Hs$, $\Delta_1=1$ y $T=G/(1+G\Hs)$, la regla \eqref{eq:so-retro}.

\begin{ejemplo}[label=ej:so-mason]{Dos trayectos y cuatro lazos}
Halle $Y(s)/R(s)$ para el grafo de la \cref{fig:so-mason}.
\solucion \textbf{Trayectos directos:} $P_1=G_1G_2G_3G_4G_5G_6$ (nodos 1-2-3-4-5-6-7) y $P_2=G_1G_2G_7G_8G_9$ (nodos 1-2-3-8-9-7).

\textbf{Lazos:} $L_1=-G_2H_1$ (nodos 2-3), $L_2=-G_2G_3G_4H_4$ (2-3-4-5), $L_3=-G_4H_2$ (4-5) y $L_4=-G_8H_3$ (8-9).

\textbf{Lazos disyuntos:} en pares, $L_1L_3$, $L_1L_4$, $L_2L_4$ y $L_3L_4$ ($L_1$ y $L_2$ comparten los nodos 2 y 3; $L_2$ y $L_3$, el 4 y el 5). En ternas, solo $L_1L_3L_4$.
\[
  \Delta=1-(L_1+L_2+L_3+L_4)+(L_1L_3+L_1L_4+L_2L_4+L_3L_4)-L_1L_3L_4.
\]
\textbf{Cofactores:} $P_1$ toca los nodos 1 a 7, así que solo queda libre $L_4$: $\Delta_1=1-L_4$. $P_2$ toca 1, 2, 3, 8, 9 y 7, así que solo queda libre $L_3$: $\Delta_2=1-L_3$. Finalmente,
\[
  \frac{Y(s)}{R(s)}=\frac{P_1(1-L_4)+P_2(1-L_3)}{\Delta}.
\]
\end{ejemplo}

\begin{figure}[ht]
\centering
\includegraphics[width=0.95\linewidth]{cap03/mason_grafo}
\caption{Ejemplo~\ref{ej:so-mason}: grafo de flujo de señal con nueve nodos, dos trayectos directos y cuatro lazos (en coral, las ramas de retroalimentación).}
\label{fig:so-mason}
\end{figure}

\codigorepo{cap03_modelado_segundo_orden/C3_Ej6_Mason.m}{mason}{ejemplo \ref{ej:so-mason}: Mason contra resolver las ecuaciones de los nodos}

\begin{ideaclave}
El discriminante es siempre el denominador de la función de transferencia: \alerta{$\Delta(s)=0$ es la ecuación característica} y sus raíces son los polos del lazo cerrado. Para el lazo simple, $\Delta=1+G\Hs$: lo mismo que encontramos con el álgebra de bloques.
\end{ideaclave}

\begin{receta}[Aplicar la ley de Mason]
\begin{enumerate}
  \item Numerar los nodos. Listar \emph{todos} los trayectos directos con sus nodos.
  \item Listar todos los lazos con sus nodos y sus ganancias (incluyendo el signo de cada retroalimentación).
  \item Buscar pares, ternas\ldots\ de lazos que no compartan nodos y armar $\Delta$.
  \item Para cada trayecto, tachar los lazos que lo tocan y armar $\Delta_k$ con los que quedan.
  \item Aplicar \eqref{eq:so-mason} y, si se puede, verificar con álgebra de bloques o con \matlab.
\end{enumerate}
\end{receta}

\begin{pausa}
En la \cref{fig:so-mason}, si se elimina la rama $-H_4$, ¿cuánto vale el nuevo $\Delta$?
\respuesta{Quedan $L_1$, $L_3$ y $L_4$, todos disyuntos entre sí: $\Delta=1-(L_1+L_3+L_4)+(L_1L_3+L_1L_4+L_3L_4)-L_1L_3L_4=(1-L_1)(1-L_3)(1-L_4)$. Cuando ningún lazo se toca, el discriminante se factoriza.}
\end{pausa}

%% ---------------------------------------------------------------------
\section{Sistemas de orden superior y polos dominantes}
\label{sec:so-ordensuperior}
%% ---------------------------------------------------------------------

Un lazo cerrado con controlador, válvula y sensor rara vez es de segundo orden. ¿Sirve entonces todo lo anterior? Sí, gracias a una idea: los \clave{polos dominantes}, los que contienen la mayor parte de la información del comportamiento. Si los demás polos están lejos (a la izquierda), sus modos mueren rápido y la respuesta se parece a la de un segundo orden. Así no hace falta antitransformar para estimar $\OS$ y $\ts$.

\paragraph{Cuándo despreciar un polo} Para $T(s)=\dfrac{1}{(s^2+2\zeta\wn s+\wn^2)(\gamma s+1)}$, P.\,R.\ Clement mostró que el $\OS$ y el $\ts$ se parecen a los del segundo orden cuando~\citep{Cle60}
\begin{equation}
  \left|\frac1\gamma\right|\geq10\,|\zeta\wn|,
  \label{eq:so-clement}
\end{equation}
es decir, cuando la parte real del polo adicional es al menos diez veces la de los dominantes. Es una regla conservadora: con cinco veces suele bastar en la práctica.

\begin{ejemplo}[label=ej:so-dominantes]{Un polo lejano}
Para $T(s)=\dfrac{20}{(s+10)(s^2+2s+2)}$ obtenga la respuesta al escalón y compárela con la del segundo orden dominante.
\solucion Los polos dominantes están en $-1\pm j$ ($\zeta\wn=1$) y el tercero en $-10$: se cumple \eqref{eq:so-clement}. Con fracciones parciales,
\[
  Y(s)=\frac1s-\frac{1/41}{s+10}-\frac{\frac{40}{41}(s+1)+\frac{50}{41}}{(s+1)^2+1},
  \qquad
  y(t)=1-\tfrac{40}{41}e^{-t}\cos t-\tfrac{50}{41}e^{-t}\sin t-\tfrac1{41}e^{-10t}.
\]
El término del polo en $-10$ es pequeño ($1/41$) \emph{y} rápido. La aproximación debe conservar la ganancia DC: $T(s)\approx\dfrac{20/10}{s^2+2s+2}$. Los sobrepicos son $4.27\%$ y $4.32\%$ (\cref{fig:so-dominantes}, izquierda).
\end{ejemplo}

\paragraph{Los ceros también cuentan} El análisis anterior supone que no hay ceros cerca de los polos dominantes. Un cero en $-a$ cerca de ellos cambia mucho la respuesta: en $T(s)=\dfrac{(\wn^2/a)(s+a)}{s^2+2\zeta\wn s+\wn^2}$ el sobrepico depende de $a/(\zeta\wn)$, y crece cuanto más cerca está el cero. Con un polo y un cero adicionales, $T(s)=\dfrac{(\wn^2/a)(s+a)}{(s^2+2\zeta\wn s+\wn^2)(\tau s+1)}$, ambos se pueden ignorar si $a\gg\zeta\wn$ y $1/\tau\gg\zeta\wn$; si no, no.

\begin{ejemplo}[label=ej:so-cero]{Un cero y un polo que no se pueden ignorar}
Analice $T(s)=\dfrac{62.5\,(s+2.5)}{(s^2+6s+25)(s+6.25)}$.
\solucion La ganancia DC es $1$. Los dominantes tienen $\zeta=0.6$, $\wn=5$ ($\zeta\wn=3$), el cero está en $-2.5$ y el polo extra en $-6.25$ ($\tau=0.16$). Ninguno está lejos. Comparemos tres modelos (\cref{fig:so-dominantes}, derecha):
\begin{center}\small
\begin{tabular}{@{}lccc@{}}
\toprule
modelo & $\OS$ & $\ts$ (s) & comentario\\
\midrule
completo & $38.0\%$ & $1.59$ & la respuesta real\\
sin el polo: $\dfrac{10(s+2.5)}{s^2+6s+25}$ & $54.5\%$ & $1.47$ & el cero solo dispara el sobrepico\\
solo los dominantes: $\dfrac{25}{s^2+6s+25}$ & $9.5\%$ & $1.19$ & la fórmula del segundo orden\\
\bottomrule
\end{tabular}
\end{center}
El tercer polo amortigua el sobrepico que provoca el cero, pero alarga el $\ts$; ninguna de las dos aproximaciones sirve.
\end{ejemplo}

\begin{figure}[ht]
\centering
\includegraphics[width=\linewidth]{cap03/f_dominantes}
\caption{Izquierda: ejemplo~\ref{ej:so-dominantes}, el polo en $-10$ se puede despreciar. Derecha: ejemplo~\ref{ej:so-cero}, el cero en $-2.5$ y el polo en $-6.25$ no.}
\label{fig:so-dominantes}
\end{figure}

\codigorepo{cap03_modelado_segundo_orden/C3_Ej7_Dominantes.m}{cero}{ejemplo \ref{ej:so-cero}}

\subsection{Ubicación de las raíces y respuesta transitoria}

En general, el lazo cerrado es $T(s)=\dfrac{\sum_kP_k\Delta_k}{\Delta(s)}$ y sus polos son las raíces de $\Delta(s)=0$ (para un lazo simple, $1+G\Hs=0$). Sin raíces repetidas y con ganancia unitaria, la respuesta al escalón tiene la forma
\[
  y(t)=1+\sum_{i}A_i\,e^{-\sigma_it}+\sum_kD_k\,e^{-\alpha_kt}\sin(\omega_kt+\theta_k),
\]
una suma de exponenciales (polos reales $-\sigma_i$) y senoides amortiguadas (polos $-\alpha_k\pm j\omega_k$). Para que la salida sea acotada todas las partes reales deben ser negativas (\cref{fig:po-modos}). Y para leer la forma de la respuesta:
\begin{itemize}
  \item La respuesta está determinada por los polos \emph{y} por los residuos.
  \item Los polos más cercanos al eje imaginario son los modos más lentos: dominan la forma de la respuesta, a menos que su residuo sea pequeño.
  \item Los polos fijan \emph{qué} modos hay; los ceros fijan \emph{cuánto} pesa cada uno. Un cero cerca de un polo reduce la contribución de ese modo (en el límite, lo cancela).
\end{itemize}

\begin{pausa}
En el ejemplo~\ref{ej:so-dominantes}, ¿cuánto pesaría el modo $e^{-10t}$ si el polo extra estuviera en $-3$ en lugar de $-10$ (manteniendo la ganancia DC)? ¿Se podría despreciar?
\respuesta{Con $T=\dfrac{6}{(s+3)(s^2+2s+2)}$, el residuo de $Y=T/s$ en $-3$ es $\dfrac{6}{(-3)(9-6+2)}=-0.4$: casi la mitad del valor final, y el modo $e^{-3t}$ es apenas tres veces más rápido que la envolvente $e^{-t}$. No se puede despreciar.}
\end{pausa}

%% ---------------------------------------------------------------------
\section{Simplificación de sistemas lineales}
\label{sec:so-reduccion}
%% ---------------------------------------------------------------------

A veces conviene \emph{reemplazar} un modelo de orden alto por uno más simple, por ejemplo para diseñar un controlador con las fórmulas del segundo orden. Veremos dos maneras.

\paragraph{1. Despreciar un polo rápido, conservando la ganancia} Si un polo real negativo tiene magnitud mucho mayor que los demás, se elimina, pero dividiendo por su valor para no cambiar la ganancia en baja frecuencia. Por ejemplo,
\[
  G(s)=\frac{K}{s(s+2)(s+30)}\approx\frac{K/30}{s(s+2)},
\]
porque para $|s|\ll30$, $s+30\approx30$. Con $K=60$ y retroalimentación unitaria los sobrepicos son $5.3\%$ y $4.3\%$ y las respuestas casi se superponen (\cref{fig:so-reduccion}, izquierda). La aproximación empeora cuando el lazo es más agresivo: con $K=300$ los polos dominantes del lazo cerrado se alejan del origen, se acercan en magnitud al de $-30$, y los sobrepicos pasan a $42.3\%$ y $35.1\%$.

\codigorepo{cap03_modelado_segundo_orden/C3_Ej8_Reduccion.m}{despreciar}{despreciar el polo en $-30$}

\paragraph{2. Igualación de momentos} Cuando ningún polo es despreciable hay que buscar un modelo de orden menor que se \emph{parezca} en frecuencia al original~\citep{DoBi08}. Escribimos el modelo de orden alto y el de orden bajo con la misma ganancia $K$ (para conservar el estado estacionario) y con término independiente $1$:
\[
  G_H(s)=K\,\frac{a_ms^m+\cdots+a_1s+1}{b_ns^n+\cdots+b_1s+1},\qquad
  G_L(s)=K\,\frac{c_ps^p+\cdots+c_1s+1}{d_gs^g+\cdots+d_1s+1},\quad p\leq g<n.
\]
Sea $\dfrac{G_H(s)}{G_L(s)}=\dfrac{M(s)}{\Delta(s)}$ (aquí $\Delta$ no es el discriminante de Mason, solo un denominador). Se definen
\begin{equation}
  M_{2q}=\sum_{k=0}^{2q}\frac{(-1)^{k+q}\,M^{(k)}(0)\,M^{(2q-k)}(0)}{k!\,(2q-k)!},\qquad q=0,1,2,\dots
  \label{eq:so-momentos}
\end{equation}
($M^{(k)}$ es la $k$-ésima derivada) y lo mismo para $\Delta_{2q}$. Los coeficientes desconocidos $c_i,d_i$ se obtienen igualando
\[
  M_{2q}=\Delta_{2q},\qquad q=1,2,\dots
\]
hasta tener tantas ecuaciones como incógnitas. (Como $M^{(k)}(0)/k!$ es el coeficiente de $s^k$, la fórmula solo usa coeficientes.) La idea: $|G_H(\jw)/G_L(\jw)|^2=\dfrac{\sum_q M_{2q}\omega^{2q}}{\sum_q\Delta_{2q}\omega^{2q}}$, así que igualar los primeros momentos hace que ese cociente se parezca a $1$ en baja frecuencia.

\begin{ejemplo}[label=ej:so-momentos]{De tercer a segundo orden}
Reduzca $G_H(s)=\dfrac{6}{s^3+6s^2+11s+6}=\dfrac{1}{\frac16s^3+s^2+\frac{11}6s+1}$ a $G_L(s)=\dfrac{1}{d_2s^2+d_1s+1}$.
\solucion $\dfrac{G_H}{G_L}=\dfrac{1+d_1s+d_2s^2}{\frac16s^3+s^2+\frac{11}6s+1}$, así que $M(s)=1+d_1s+d_2s^2$ y $\Delta(s)=\frac16s^3+s^2+\frac{11}6s+1$. Las derivadas en cero son $M^{(0)}=1$, $M^{(1)}=d_1$, $M^{(2)}=2d_2$; $\Delta^{(0)}=1$, $\Delta^{(1)}=\frac{11}6$, $\Delta^{(2)}=2$, $\Delta^{(3)}=1$. Con dos incógnitas usamos $q=1,2$:
\[
  M_2=d_1^2-2d_2,\quad \Delta_2=\tfrac{121}{36}-2=\tfrac{49}{36};\qquad
  M_4=d_2^2,\quad \Delta_4=1-2\cdot\tfrac{11}{6}\cdot\tfrac16=\tfrac{7}{18}.
\]
Entonces $d_2=\sqrt{7/18}=0.6236$ y $d_1=\sqrt{49/36+2d_2}=1.615$ (se toman las raíces positivas: con las negativas $G_L$ tendría polos inestables). Así
\[
  G_L(s)=\frac{1}{0.6236s^2+1.615s+1}=\frac{1.604}{s^2+2.590s+1.604},
\]
con polos en $-1.024$ y $-1.565$ (los de $G_H$ estaban en $-1$, $-2$ y $-3$). La \cref{fig:so-reduccion} (derecha) compara los escalones: la diferencia máxima es de $0.10$, al principio, porque $G_L$ tiene grado relativo $2$ y arranca más rápido.
\end{ejemplo}

\begin{figure}[ht]
\centering
\includegraphics[width=\linewidth]{cap03/f_reduccion}
\caption{Izquierda: despreciar el polo en $-30$ (lazo cerrado, $K=60$). Derecha: ejemplo~\ref{ej:so-momentos}, igualación de momentos.}
\label{fig:so-reduccion}
\end{figure}

\codigorepo{cap03_modelado_segundo_orden/C3_Ej8_Reduccion.m}{momentos}{ejemplo \ref{ej:so-momentos} con \cmd{reduccionMomentos}}

%% ---------------------------------------------------------------------
\section{Para llevar}
\label{sec:so-cierre}
%% ---------------------------------------------------------------------

\begin{parallevar}
\begin{enumerate}
  \item Linealizar es reemplazar cada no linealidad por su recta tangente en el punto de operación: $\dot{\dev{x}}=\sum\frac{\partial f}{\partial x_j}\big|_\ast\dev{x}_j+\sum\frac{\partial f}{\partial u_k}\big|_\ast\dev{u}_k$. Solo vale cerca del punto.
  \item Segundo orden: $K\wn^2/(s^2+2\zeta\wn s+\wn^2)$. $\zeta$ decide si oscila; $\wn$, qué tan rápido.
  \item Subamortiguado: $\tp=\pi/\wdamp$ (parte imaginaria), $\OS=100e^{-\zeta\pi/\sqrt{1-\zeta^2}}$ (ángulo), $\ts\approx4/(\zeta\wn)$ (parte real), $\tr=(\pi-\cos^{-1}\zeta)/\wdamp$ ($0$--$100\%$). Especificaciones de tiempo $\leftrightarrow$ región del plano-$s$.
  \item Bloques: serie multiplica, paralelo suma, retroalimentación da $G/(1+G\Hs)$. $1+G\Hs=0$ es la ecuación característica; $G\Hs$ es la malla abierta.
  \item Mason: $T=\frac1\Delta\sum P_k\Delta_k$. El discriminante $\Delta$ es el denominador: $\Delta=0$ es la ecuación característica.
  \item Polos dominantes: un polo extra se desprecia si está unas diez veces más a la izquierda (Clement), \emph{conservando la ganancia DC}; un cero cerca de los dominantes cambia mucho el sobrepico.
  \item Reducir orden: despreciar polos rápidos o igualar momentos. Siempre verificar con una simulación.
\end{enumerate}
\end{parallevar}

\begin{profundiza}[Correcciones respecto a la versión anterior del capítulo]
(i) El ejemplo numérico de los dos tanques ($\tau_1=1$, $\tau_2=4$, $C_2=2$) no era físico: en esa planta $C_2=-C_1$ y $\zeta\geq1$ siempre. Se reemplazó por datos físicos y el subamortiguado se ilustra con la masa-resorte. (ii) $\beta=\cos^{-1}\zeta$ (decía $\beta=\cos\zeta$). (iii) La condición para polos complejos de los tanques estaba mal escrita; con $C_2=-C_1$ nunca se cumple. (iv) El sobrepico se escribe $\OS$ (antes $MP$, que se confundía con el pico de resonancia), ``tiempo de establecimiento'' (no ``de asentamiento'') y $\Delta$ es el \emph{discriminante}. (v) Ejemplo con cero y polo adicionales: $54.5\%$ y $38.0\%$ (antes $55\%$ y $38\%$) y $\ts=1.47$ y $1.59$\,s (antes $1.33$, que era la fórmula $4/(\zeta\wn)$, y $1.6$). (vi) La aproximación de $\tr$ de $10$--$90\%$ se verificó: error menor que el $3.1\%$ para $0.1\leq\zeta\leq0.9$. (vii) El ejercicio de reducción de orden con un cero tenía una respuesta que no satisface las ecuaciones de momentos; se reformuló. Todo se verificó con \texttt{herramientas/cap03\_numeros.py} y \cmd{verificarCap03.m}.
\end{profundiza}

%% ---------------------------------------------------------------------
\section{Problemas}
\label{sec:so-problemas}
%% ---------------------------------------------------------------------

Los problemas marcados con \etiquetaM{} requieren \matlab/\simulink; los marcados con \etiquetaEstrella{} son de profundización.

\subsection*{Linealización}

\begin{problema}[linealización]\label{prob:so-lin}
Linealice alrededor del equilibrio y obtenga la función de transferencia, la ganancia y la constante de tiempo: (a) $\dot x=-x+2\sqrt x+u$ con $\eq{u}=0$ y $\eq{x}>0$; (b) $\dot x=-x^3+2xu$ con $\eq{u}=2$ (equilibrio no trivial). (c) \etiquetaM{} Para el tanque de la pausa de la \cref{sec:so-linealizacion}, grafique $c\sqrt h$ y su aproximación lineal y encuentre el rango de alturas en el que el error relativo es menor que el $5\%$. Compárelo con la \cref{fig:so-lin1}.
\end{problema}

\begin{problema}[estructura de los tanques]\label{prob:so-bloques-tanques}
Dibuje el diagrama de bloques que corresponde a las dos ecuaciones de $\Dev{H}_1$ y $\Dev{H}_2$ de la \cref{sec:so-tanques} (un lazo de retroalimentación con ganancia $C_2$). Redúzcalo con \eqref{eq:so-retro} y con la ley de Mason y verifique \eqref{eq:so-tanques-ft}. ¿De qué signo es la retroalimentación cuando $C_2=-C_1$?
\end{problema}

\begin{problema}[tanques con otros datos]\label{prob:so-tanques}
Con $A_1=1$\,m$^2$, $A_2=2$\,m$^2$, $C'_{v1}=0.2$, $C'_{v2}=0.1$\,m$^{5/2}$/s y $\eq{h_2}=1$\,m: (a) halle $\eq{f_i}$ y $\eq{h_1}$; (b) $C_1$, $C_2$, $C_3$, $K_1$, $\tau_1$, $K_2$, $\tau_2$; (c) la función de transferencia, su ganancia, $\wn$, $\zeta$ y sus polos; (d) la respuesta ante un escalón de $0.01$\,m$^3$/s y su $\ts$; (e) verifique que en el nuevo estado estacionario $h_1>h_2$ (si no, el modelo \eqref{eq:so-valvulas} pierde sentido); (f) \etiquetaM{} adapte \cmd{crearModelo\_Tanques.m} y compare los modelos lineal y no lineal con escalones de $0.01$ y $0.05$\,m$^3$/s.
\end{problema}

\begin{problema}[los tanques nunca oscilan]\label{prob:so-zeta-tanques}
\etiquetaEstrella{} Con $A_1,A_2,C'_{v1},C'_{v2}>0$ arbitrarios: (a) muestre que $C_2=-C_1$ y $C_1,C_3>0$ en cualquier punto de operación admisible; (b) muestre que $\wn^2=\dfrac{C_1C_3}{A_1A_2}$ y $2\zeta\wn=\dfrac{C_1+C_3}{A_2}+\dfrac{C_1}{A_1}$; (c) con la desigualdad de las medias aritmética y geométrica, demuestre que $\zeta\geq1$. ¿Cuándo se alcanzaría $\zeta=1$? ¿Es realizable? (d) Explique por qué el ejemplo con $C_2=2$ de la versión anterior no corresponde a ningún par de tanques. (e) Proponga una modificación física (por ejemplo, una bomba en lugar de la válvula entre tanques) que permita $C_2>0$ y diga cuándo el sistema sería subamortiguado.
\end{problema}

\subsection*{Segundo orden y diseño}

\begin{problema}[cuatro sistemas]\label{prob:so-cuatro}
Para $T_1=\dfrac{50}{s^2+6s+25}$, $T_2=\dfrac{16}{s^2+16}$, $T_3=\dfrac{9}{s^2+6s+9}$ y $T_4=\dfrac{100}{s^2+30s+100}$: (a) identifique $K$, $\wn$, $\zeta$, clasifíquelos y ubique sus polos; (b) para el subamortiguado calcule $\tp$, $\OS$, $\ts$ al $2$ y al $5\%$, y $\tr$ con las dos definiciones; (c) para el crítico, con \eqref{eq:so-critico}, halle el instante en que llega al $95\%$; (d) para el sobreamortiguado, construya la aproximación $s_2/(s+s_2)$ y cuantifique el error máximo ante el escalón; (e) \etiquetaM{} verifique con \cmd{stepinfo} y explique las diferencias en $\tr$.
\end{problema}

\begin{problema}[diseño con tacómetro]\label{prob:so-diseno}
En la estructura de la \cref{fig:so-kkh}, con camino directo $\dfrac{K}{s(s+4)}$: (a) muestre que $\dfrac YR=\dfrac{K}{s^2+(4+KK_h)s+K}$; (b) halle $K$ y $K_h$ para $\OS=10\%$ y $\ts=1.5$\,s; (c) calcule $\tp$, $\tr$ y los polos; (d) si además se pide $\tp\leq0.8$\,s, ¿es compatible? Si no, ¿qué $\ts$ lo permitiría con el mismo $\OS$? Explíquelo con la \cref{fig:so-lineas}; (e) \etiquetaM{} verifique con \cmd{series} y \cmd{feedback}.
\end{problema}

\begin{problema}[región de diseño]\label{prob:so-region}
Dibuje en el plano-$s$ la región donde deben quedar los polos de un segundo orden para que $\OS\leq10\%$ y $\ts\leq2$\,s. ¿Cumplen los polos $-3\pm j3$? ¿Y $-1\pm j1$?
\end{problema}

\subsection*{Bloques y Mason}

\begin{problema}[sistema multilazo]\label{prob:so-multilazo}
Las señales cumplen $E=R-H_2Y$, $V=G_1E$, $W=G_2(V-H_1W)$, $Y=G_3W$, con $G_1=\frac1{s+1}$, $G_2=\frac2s$, $G_3=\frac5{s+3}$, $H_1=1$ y $H_2=0.5$. (a) Dibuje el diagrama de bloques. (b) Obtenga $T=Y/R$ reduciendo primero el lazo interno. (c) Obtenga $T$ con Mason (trayectos, lazos, disyuntos, $\Delta$ y $\Delta_k$). (d) Escriba la ecuación característica y sus raíces. (e) \etiquetaM{} Verifique con \cmd{feedback}.
\end{problema}

\subsection*{Orden superior y reducción}

\begin{problema}[tres aproximaciones]\label{prob:so-tres}
Para $T(s)=\dfrac{10}{s^3+6s^2+11s+11}$ (el del problema anterior): (a) halle los polos y verifique el criterio \eqref{eq:so-clement}; (b) construya la aproximación por polos dominantes (conservando la ganancia DC) y prediga $\OS$ y $\ts$; (c) obtenga otra por igualación de momentos; (d) \etiquetaM{} simule las tres, compare $\OS$ y $\ts$ y decida cuál es preferible; (e) explique por qué el criterio de (a), aunque no se cumpla, no impidió una aproximación aceptable.
\end{problema}

\begin{problema}[momentos con un cero]\label{prob:so-momentos}
\etiquetaM{} Reduzca $G_H(s)=\dfrac{s^3+7s^2+24s+24}{(s+1)(s+2)(s+3)(s+4)}$ a $G_L(s)=\dfrac{1}{d_2s^2+d_1s+1}$ por igualación de momentos (note que ahora $M(s)$ incluye el numerador de $G_H$). Compare los escalones. \etiquetaEstrella{} Intente luego $G_L=\dfrac{1+c_1s}{d_2s^2+d_1s+1}$ con $q=1,2,3$: ¿tiene solución real?
\end{problema}

\begin{problema}[un resorte no lineal]\label{prob:so-resorte}
\etiquetaEstrella{} Una masa con resorte endurecido y amortiguamiento cuadrático cumple $m\ddot x+b\,\dot x|\dot x|+k_1x+k_3x^3=F$, con $m=1$\,kg, $b=0.8$\,N\,s$^2$/m$^2$, $k_1=4$\,N/m y $k_3=1$\,N/m$^3$. (a) Halle el equilibrio para $\eq{F}=5$\,N. (b) Linealice y obtenga $\Dev{X}/\Dev{F}$, $\wn$ y $\zeta$. (c) El $\zeta$ debe llamarle la atención: explique por qué el amortiguamiento desaparece al linealizar y qué predice el modelo lineal. (d) \etiquetaM{} Simule el no lineal con un escalón adicional de $0.5$\,N y compárelo con el lineal; defina un criterio y diga durante cuántos ciclos es aceptable la predicción lineal. (e) Estime un $\zeta$ equivalente con el decremento logarítmico. (f) Repita (a)--(b) con $\eq{F}=16$\,N y explique por qué un solo modelo LTI no describe todo el rango.
\end{problema}

%% ---------------------------------------------------------------------
\section{Soluciones seleccionadas}
%% ---------------------------------------------------------------------

\paragraph{Problema~\ref{prob:so-lin}}
(a) $\eq{x}=4$; $\partial f/\partial x=-1+1/\sqrt{\eq{x}}=-0.5$, $\partial f/\partial u=1$: $\dfrac{\Dev{X}}{\Dev{U}}=\dfrac{1}{s+0.5}=\dfrac{2}{2s+1}$ ($K_1=2$, $\tau=2$). (b) $\eq{x}=2$; $-3(\eq{x})^2+2\eq{u}=-8$ y $2\eq{x}=4$: $\dfrac{4}{s+8}=\dfrac{0.5}{0.125s+1}$. (c) $2.13\leq h\leq7.51$\,m: un rango asimétrico alrededor de $\eq{h}=4$, porque la raíz cuadrada se curva más para alturas bajas.

\paragraph{Problema~\ref{prob:so-bloques-tanques}}
Es un lazo con camino directo $\frac{K_1}{\tau_1s+1}\frac{K_2}{\tau_2s+1}$ y retroalimentación $C_2$ (signo $-$ en el comparador): $T=\dfrac{K_1K_2}{(\tau_1s+1)(\tau_2s+1)+K_1K_2C_2}$. Con $C_2=-C_1<0$ la retroalimentación es \emph{positiva}: al subir $h_2$ disminuye $f_1$, el tanque 1 se llena más y empuja más agua hacia el 2. Pero su ganancia de malla, $K_1K_2C_1=C_1/(C_1+C_3)<1$, no alcanza a inestabilizar el sistema.

\paragraph{Problema~\ref{prob:so-tanques}}
(a) $\eq{f}=0.1$\,m$^3$/s, $\eq{h_1}=1.25$\,m. (b) $C_1=0.2$, $C_2=-0.2$, $C_3=0.05$\,m$^2$/s; $K_1=5$, $\tau_1=5$\,s, $K_2=0.8$, $\tau_2=8$\,s. (c) $\dfrac{20}{200s^2+65s+1}$: ganancia $20$\,s/m$^2$, $\wn=0.0707$\,rad/s, $\zeta=2.30$, polos $-0.0162$ y $-0.309$. (d) $\Delta h_2\to0.2$\,m con $\ts=245$\,s (la estimación con el polo lento, $4/0.0162=247$\,s). (e) $h_1=1.25+0.05+0.2=1.5>h_2=1.2$\,m.

\paragraph{Problema~\ref{prob:so-zeta-tanques}}
(c) $2\zeta\wn=\dfrac{C_1}{A_2}+\dfrac{C_3}{A_2}+\dfrac{C_1}{A_1}\geq\dfrac{C_3}{A_2}+\dfrac{C_1}{A_1}\geq2\sqrt{\dfrac{C_1C_3}{A_1A_2}}=2\wn$, luego $\zeta\geq1$. La primera desigualdad es estricta porque $C_1/A_2>0$: $\zeta=1$ exigiría $C_1\to0$ (válvula entre tanques cerrada), que desacopla los tanques.

\paragraph{Problema~\ref{prob:so-cuatro}}
(a) $T_1$: $K=2$, $\wn=5$, $\zeta=0.6$; $T_2$: $K=1$, $\wn=4$, $\zeta=0$; $T_3$: $K=1$, $\wn=3$, $\zeta=1$; $T_4$: $K=1$, $\wn=10$, $\zeta=1.5$ (polos $-3.82$ y $-26.2$). (b) $\tp=0.785$\,s, $\OS=9.48\%$, $\ts\approx1.33$\,s ($2\%$) y $1.0$\,s ($5\%$), $\tr=0.554$\,s ($0$--$100\%$) y $\approx0.38$\,s ($10$--$90\%$; exacto $0.371$). (c) $1.58$\,s. (d) $\dfrac{3.82}{s+3.82}$; error máximo $0.105$, al principio.

\paragraph{Problema~\ref{prob:so-diseno}}
(b) $\zeta=0.591$, $\zeta\wn=4/1.5=2.67$, $\wn=4.51$\,rad/s; $K=20.35$ y $K_h=(2\zeta\wn-4)/K=0.0655$. (c) $\tp=0.864$\,s. (d) No: $\tp\leq0.8$ con $\zeta=0.591$ exige $\wn\geq4.87$, y entonces $\ts\leq1.39$\,s: hay que mover los polos \emph{por la recta de $\zeta$ constante}, alejándolos del origen.

\paragraph{Problema~\ref{prob:so-region}}
$\zeta\geq0.591$ (un cono de semiángulo $\cos^{-1}0.591=53.8\grados$ alrededor del semieje real negativo) y $\zeta\wn\geq2$ (a la izquierda de $\Re s=-2$). $-3\pm j3$: $\zeta=0.707$ y parte real $-3$, cumple. $-1\pm j1$: $\zeta=0.707$ pero parte real $-1$, no cumple el $\ts$.

\paragraph{Problema~\ref{prob:so-multilazo}}
Lazo interno: $\frac{2/s}{1+2/s}=\frac{2}{s+2}$. $T=\dfrac{10}{(s+1)(s+2)(s+3)+5}=\dfrac{10}{s^3+6s^2+11s+11}$. Con Mason: un trayecto $P_1=G_1G_2G_3$; lazos $L_1=-G_2H_1$ y $L_2=-G_1G_2G_3H_2$, que se tocan; $\Delta=1-L_1-L_2$, $\Delta_1=1$. Polos: $-3.90$ y $-1.05\pm j1.31$.

\paragraph{Problema~\ref{prob:so-tres}}
(a) La relación entre partes reales es $3.90/1.05=3.7<10$. (b) $\dfrac{(10/11)\,2.817}{s^2+2.096s+2.817}$: $\zeta=0.624$, $\wn=1.68$; $\OS=8.1\%$ y $\ts\approx3.6$\,s. (c) $d_1=0.768$, $d_2=0.340$. (d) Original: $\OS=7.2\%$, $\ts=3.84$\,s; dominantes: $8.1\%$ y $3.56$\,s; momentos: $6.4\%$ y $3.51$\,s. (e) El residuo del polo real es pequeño.

\paragraph{Problema~\ref{prob:so-momentos}}
$d_1=1.211$ y $d_2=0.2295$ (polos $-1.03$ y $-4.25$); diferencia máxima entre escalones $\approx0.10$. Con un cero, las tres ecuaciones no tienen solución real exacta.

\paragraph{Problema~\ref{prob:so-resorte}}
(a) $4x+x^3=5\Rightarrow\eq{x}=1$\,m. (b) $k_{\text{ef}}=k_1+3k_3(\eq{x})^2=7$ y la derivada de $b\dot x|\dot x|$ en $\dot x=0$ es cero: $\dfrac{1}{s^2+7}$, $\wn=2.65$\,rad/s, $\zeta=0$. (c) El amortiguamiento cuadrático es de segundo orden en $\dot x$: la linealización lo borra y predice una oscilación eterna. (f) $\eq{x}=2$\,m, $k_{\text{ef}}=16$, $\wn=4$\,rad/s.

\respuestaspausas

